Let $a,b>1$ be coprime integers. Every sufficiently large integer
is a sum of distinct elements of
$G=\{a^u b^v:u,v\geq0\}$. We reconstruct Birch's proof through
bounded gaps, disjoint signed identities, and residue adjustment.
The intended nontrivial-base assumption matters: if one base is
one, the set consists of powers of the other base and is complete
only when that other base is two. For a base $c\geq3$, numbers
$2c^m$ cannot be sums of distinct powers of $c$.
